% 阋神星（综述）
% license CCBYSA3
% type Wiki

本文根据 CC-BY-SA 协议转载翻译自维基百科\href{https://en.wikipedia.org/wiki/Eris_(dwarf_planet)}{相关文章}。

\begin{figure}[ht]
\centering
\includegraphics[width=6cm]{./figures/88a62f7999189f0d.png}
\caption{Hubble Space Telescope 于 2006 年 8 月拍摄的 Eris 和 Dysnomia 的低分辨率图像} \label{fig_Eris_1}
\end{figure}
Eris（小行星编号：136199 Eris）是太阳系中已知质量最大、体积第二大的矮行星\(^{[23]}\)。它是一颗处于散射盘中的跨海王星天体（TNO），具有高偏心率轨道。Eris 于 2005 年 1 月由 Mike Brown 领导的 Palomar Observatory 团队发现，并于同年稍后得到确认。该天体于 2006 年 9 月以希腊—罗马神话中的纷争与不和女神命名。Eris 是围绕太阳运行的已知天体中质量第九大者，在整个太阳系中（包括卫星）质量排名第十六。它也是太阳系中已知尚未被航天器造访的最大天体。Eris 的直径测量为 2,326 ± 12 km（1,445 ± 7 mi）\(^{[13]}\)；其质量为 Earth 的 0.28\%，比 Pluto 大 27\%\(^{[24][25]}\)，尽管 Pluto 在体积上略大一些\(^{[26]}\)。Eris 与 Pluto 的表面积均可与 Russia 或 South America 相当。

Eris 已知有一颗大型卫星，名为 Dysnomia。2016 年 2 月，Eris 与太阳的距离为 96.3 AU（144.1 亿 km；89.5 亿 mi）\(^{[21]}\)，超过 Neptune 或 Pluto 距离的三倍。除长周期彗星外，直到 2018 年发现 2018 AG37 和 2018 VG18 之前，Eris 与 Dysnomia 一直是太阳系中已知距离最遥远的天然天体\(^{[21]}\)。

由于 Eris 起初看起来比 Pluto 更大，NASA 最初将其称为太阳系的第十颗行星。这一认识及未来可能发现与之大小相近的其他天体的前景，促使国际天文学联合会（IAU）首次定义“行星”一词。根据 IAU 于 2006 年 8 月 24 日批准的定义，Eris、Pluto 和 Ceres 均为“矮行星”\(^{[27]}\)，使太阳系已知行星的数量减少为八颗，与 1930 年 Pluto 被发现之前相同。2010 年一次 Eris 对恒星的掩食观测表明，它略小于 Pluto\(^{[28][29]}\)，而 New Horizons 于 2015 年 7 月测得 Pluto 的平均直径为 2,377 ± 4 km（1,477 ± 2 mi）\(^{[30][31]}\)。
\subsection{发现}  
Eris 由 Mike Brown、Chad Trujillo 和 David Rabinowitz\(^{[2]}\) 于 2005 年 1 月 5 日发现，所使用的图像拍摄于 2003 年 10 月 21 日\(^{[32]}\)。该发现于 2005 年 7 月 29 日公布，与 Makemake 的公布在同一天，且比 Haumea 晚两天\(^{[33]}\)，部分原因与后来有关 Haumea 的争议事件有关。该搜寻团队多年来一直在系统性地扫描太阳系外层的大型天体，并参与发现了数颗其他大型 TNO，包括矮行星 Quaoar、Orcus 和 Sedna\(^{[34]}\)。

该团队于 2003 年 10 月 21 日使用位于 California 的 Palomar Observatory 的 1.2 m Samuel Oschin Schmidt 望远镜进行了常规观测，但当时未能识别出 Eris 的图像，因为它在天空中的运动非常缓慢：团队的自动图像搜索软件会排除所有运动速度小于每小时 1.5 角秒的天体，以减少返回的误报数量\(^{[32]}\)。2003 年 Sedna 被发现时，其运动速度为每小时 1.75 角秒，因此团队以更低的角运动阈值重新分析了旧数据，并人工筛选了此前被排除的图像。2005 年 1 月，重新分析揭示了 Eris 相对于背景恒星的缓慢轨道运动\(^{[32]}\)。

随后进行了后续观测，以便初步确定 Eris 的轨道，从而估计该天体的距离\(^{[32]}\)。团队原计划在进一步观测和计算完成之前，暂缓公布明亮天体 Eris 和 Makemake 的发现，但在团队一直追踪的另一颗大型 TNO——Haumea——于 7 月 27 日被一个西班牙团队有争议地公布后，两个发现均于 7 月 29 日宣布\(^{[2]}\)。

Eris 的回溯预发现图像已被追溯至 1954 年 9 月 3 日\(^{[6]}\)。

2005 年 10 月公布的进一步观测表明，Eris 有一颗卫星，后来被命名为 Dysnomia。对 Dysnomia 轨道的观测使科学家能够确定 Eris 的质量，并在 2007 年 6 月计算为 \((1.66\pm0.02)\times10^{22}\) kg\(^{[24]}\)，比 Pluto 大 27\%±2\%。
\subsection{名称}  
Eris 的名称来源于希腊女神 Eris（希腊语 Ἔρις），她是纷争与不和的拟人化形象\(^{[35]}\)。该名称由 Caltech 团队于 2006 年 9 月 6 日提出，并于 2006 年 9 月 13 日正式获准\(^{[36]}\)，这一命名在此之前经历了异常漫长的阶段，在这期间该天体一直以临时编号 2003 UB313 被称呼，此编号是 IAU 根据小行星命名规则自动授予的。

名称 Eris 存在两种竞争性读法，分别使用“长元音”或“短元音”，类似单词 era 的两种不同读法\(^{[4]}\)。在英语中较为常见的形式（例如 Brown 及其学生使用的读法）是 /ˈɛrɪs/，为双音节、短元音读法\(^{[3]}\)。然而，该女神的传统英语读法为 /ˈɪərɪs/，带有长元音\(^{[5]}\)。

该名称在希腊语和拉丁语中的斜格词干为 Erid-\(^{[38]}\)，如意大利语 Eride 和俄语 Эрида（Erida）所示，因此英语形容词形式为 Eridian /ɛˈrɪdiən/\(^{[9][10]}\)。
\subsubsection{Xena}  
由于当时不确定该天体被归类为行星还是小行星（因为这两类天体适用不同的命名程序）\(^{[39]}\)，因此必须等到 2006 年 8 月 24 日 IAU 做出决定后才能正式命名\(^{[40]}\)。在此期间，该天体在公众中一度被称为 Xena。“Xena” 是发现团队内部使用的非正式昵称，灵感来自电视剧《战士公主：西娜》（Xena: Warrior Princess）的主角。据报道，发现团队一直保留“Xena”这一昵称，用于命名第一颗被发现且体积大于 Pluto 的天体。据 Brown 所说：

“我们选择它是因为它以 X 开头（即行星‘X’），听起来很神话……同时我们也希望更多女性神祇的名字能被使用（如 Sedna）。此外，当时该电视剧仍在播出，这说明我们已经搜索了相当长时间！”\(^{[41]}\)

Brown 在一次采访中表示命名过程曾一度陷入停滞：

“一位记者（Ken Chang）\(^{[42]}\) 给我打电话，他当时在《纽约时报》工作，碰巧是我大学同学……他问我：‘你们提议的名字是什么？’我说：‘这个我不能告诉你。’然后他又问：‘那你们私下互相交流时叫它什么？’……据我所知，那是我唯一一次在媒体上说出这个名字，结果传播开来，我对此只有一点点愧疚；其实我还挺喜欢这个名字。”\(^{[43]}\)
\subsubsection{选择正式名称}
\begin{figure}[ht]
\centering
\includegraphics[width=6cm]{./figures/2efad5f955645faf.png}
\caption{动画展示了用于发现 Eris 的图像中其运动情况。箭头标示出 Eris 的位置。三帧图像拍摄间隔为三个小时。} \label{fig_Eris_3}
\end{figure}
根据科普作家 Govert Schilling 的说法，Brown 最初想将该天体命名为 “Lila”，源自印度教神话中的一个概念，用以描述宇宙是 Brahman 游戏的结果\(^{[34]}\)。这一名字的发音与 Brown 新出生女儿 “Lilah” 的名字相同。Brown 注意到在正式批准之前不应公开名称，他曾在前一年为 Sedna 这么做过，并遭到强烈批评。然而，除了违反命名程序外，Sedna 这一命名并未受到其他异议，也没有出现竞争名称\(^{[44]}\)。

他在宣布发现的个人网页地址中使用了 /~mbrown/planetlila，但在围绕 Haumea 发现引发的争议混乱中忘记修改。为了避免无谓激怒更多天文学同行，他声称网页是以女儿的名字命名的，并放弃了 “Lila” 作为候选名称\(^{[34]}\)。

Brown 还曾考虑过用冥王星之妻 Persephone 作为该天体的命名\(^{[2]}\)。该名字曾多次用于科幻作品中的行星\(^{[45]}\)，并在公众中广受欢迎，在《New Scientist》杂志开展的投票中轻松获胜\(^{[46]}\)。（“Xena” 虽只是昵称，但排名第四。）然而这一选择并不可行，因为已经有一颗小行星使用该名字——399 Persephone，而 Eris 正作为小行星命名\(^{[2]}\)。

发现团队于 2006 年 9 月 6 日提议使用 Eris。Brown 决定，由于该天体长期被视为行星，因此它理应如其他行星一样使用希腊或罗马神话名称。绝大多数希腊—罗马名称已被小行星使用。Brown 将 Eris 描述为他最喜欢的女神，幸好该名字尚未被使用\(^{[43]}\)。“Eris 通过在人们之间引发争吵来制造纷争与不和”，Brown 在 2006 年说，“这颗天体也确实造成了同样的效果！”\(^{[47]}\) 该名称于 2006 年 9 月 13 日获得 IAU 批准\(^{[48][49]}\)。

尽管天文学界通常不鼓励使用行星符号，但 NASA 曾为 Eris 使用 “Eris 之手” ⟨⯰⟩（U+2BF0）符号\(^{[50]}\)。该符号源自 Discordianism，这是一种与女神 Eris 相关的宗教\(^{[51]}\)。莫斯科国立大学 Sternberg Astronomical Institute 曾使用符号 Ⓚ（U+24C0），推测要么是《普林西皮亚·纷争论》中 “权利反转” 符号，要么是刻有希腊语 “Kallisti”（意为 “献给最美者”）的纷争金苹果的简化形式，该符号曾在 Discordian 讨论板上被建议作为矮行星的符号，最终定为 ⟨⯰⟩\(^{[52][53]}\)。多数占星师使用 “Eris 之手”，但偶尔也能看到其他符号，如 ⟨⯱⟩（U+2BF1）\(^{[51]}\)。
\subsection{分类}  
Eris 是一颗跨海王星矮行星。其轨道特征更具体地将其分类为散射盘天体（SDO），即一颗在太阳系形成过程中与 Neptune 发生引力相互作用后从 Kuiper 带 “散射” 到更远、更独特轨道的 TNO。尽管其轨道倾角在已知 SDO 中显得不同寻常，理论模型显示，最初位于 Kuiper 带内缘附近的天体被散射到较高倾角轨道，而来自外缘的天体则倾角较小\(^{[54]}\)。

由于 Eris 最初被认为比 Pluto 更大，NASA 与媒体在相关报道中将其描述为 “太阳系第十颗行星”\(^{[55]}\)。为回应其状态的不确定性，并因为关于 Pluto 是否应被归类为行星的争论持续存在，IAU 委托一组天文学家制定一个足够精确的 “行星” 定义来解决该问题。该定义作为 IAU 的《太阳系行星定义》于 2006 年 8 月 24 日通过。此时，Eris 与 Pluto 均被归类为矮行星，与新的行星定义有所区别\(^{[56]}\)。Brown 此后表示支持这一分类\(^{[57]}\)。IAU 随后将 Eris 加入小行星目录，并将其编号为 (136199) Eris\(^{[40]}\)。
\subsection{分类}
\begin{figure}[ht]
\centering
\includegraphics[width=10cm]{./figures/cb855c52f457014f.png}
\caption{Eris（蓝色）轨道与 Saturn、Uranus、Neptune 和 Pluto（白色/灰色）轨道的比较。黄道面以下的轨道弧段以较深颜色绘制，红点表示 Sun。左侧示意图为极视图，右侧示意图为来自黄道平面的不同视角。} \label{fig_Eris_4}
\end{figure}
\begin{figure}[ht]
\centering
\includegraphics[width=10cm]{./figures/37f9f417926c5918.png}
\caption{从 Earth 上观察，Eris 在 Cetus 星座中呈现出微小的回环运动。} \label{fig_Eris_5}
\end{figure}
Eris 的轨道周期为 561 年\(^{[6]}\)。其与 Sun 的最大可能距离（远日点）为 97.7 AU，最近距离（近日点）为 38.4 AU\(^{[6]}\)。由于近日点时间是使用未受扰动的二体解在所选历元下定义的，因此历元越接近日点日期，结果越准确；要精确预测近日点时间需要进行数值积分。JPL Horizons 的数值积分显示，Eris 在约 1699 年经过近日点\(^{[58]}\)，在约 1977 年达到远日点，并将在 2257 年 12 月前后再次回到近日点\(^{[12]}\)。不同于八大行星，它们的轨道大致都位于与 Earth 相同的平面上，Eris 的轨道倾角很大：轨道相对黄道面倾斜约 44°\(^{[6]}\)。在被发现时，Eris 及其卫星是太阳系中已知距离最遥远的天体（不包括长周期彗星和空间探测器）\(^{[2][59]}\)。这一身份一直保持到 2018 年发现 2018 VG18 为止\(^{[60]}\)。

Eris 的轨道偏心率很高，使其能够接近 Sun 到 37.9 AU，这是典型的散射盘天体近日点距离\(^{[61]}\)。这一距离位于 Pluto 轨道范围内，但仍足以避免与 Neptune 发生直接相互作用（约 37 AU）\(^{[62]}\)。而 Pluto 则不同，作为 plutinos 之一，其轨道倾角和偏心率较低，并在轨道共振保护下可以穿过 Neptune 的轨道\(^{[63]}\)。在约 800 年后，Eris 将在一段时间内比 Pluto 更接近 Sun（参见左侧图表）。

截至 2007 年，Eris 的视星等为 18.7，使其亮度足以被一些业余望远镜观测到\(^{[64]}\)。在良好条件下，配备 CCD 的 200 mm（7.9 in）望远镜即可探测到 Eris\(^{[f]}\)。之所以至今才被发现，是因为其轨道倾角很大，而对太阳系外层大型天体的搜寻通常集中在大多数天体分布的黄道面附近\(^{[65]}\)。

由于其轨道倾角很大，Eris 仅穿过传统黄道带的少数几个星座；目前它位于 Cetus 星座。它从 1876 年到 1929 年位于 Sculptor 星座，从约 1840 年到 1875 年位于 Phoenix 星座。它将在 2036 年进入 Pisces，并停留到 2065 年，届时将进入 Aries\(^{[66]}\)。随后它将移入北天，在 2128 年进入 Perseus，并在 2173 年进入 Camelopardalis（在此处达到其最北赤纬）。
\subsection{大小、质量与密度}
\begin{figure}[ht]
\centering
\includegraphics[width=6cm]{./figures/5129c0181c92affe.png}
\caption{尺寸比较：Eris（左下）与 Moon 和 Earth（上方与右侧）} \label{fig_Eris_6}
\end{figure}
\begin{figure}[ht]
\centering
\includegraphics[width=6cm]{./figures/de276f374324d3f5.png}
\caption{2010 年 11 月 Eris 对恒星掩食的示意图。由掩食弦得到的投影呈现出 Eris 的圆形轮廓，其球形直径为 2,326 km（1,445 mi）} \label{fig_Eris_7}
\end{figure}
\begin{figure}[ht]
\centering
\includegraphics[width=6cm]{./figures/174b1523ffaa1cfd.png}
\caption{} \label{fig_Eris_8}
\end{figure}
2010 年 11 月，Eris 成为迄今为止从 Earth 观测到的最遥远恒星掩食事件之一\(^{[14]}\)。该事件的初步数据对此前的大小估计提出质疑\(^{[14]}\)。研究团队于 2011 年 10 月公布掩食结果，估计其直径为 2326±12 km\(^{[13]}\)。

这使得 Eris 比直径为 2376.6±1.6 km 的 Pluto 稍小，尽管 Eris 的质量更大。结果还显示其几何反照率为 0.96。普遍推测高反照率来源于其轨道偏心导致 Eris 在接近与远离 Sun 之间温度波动，从而不断补充表面的冰层\(^{[22]}\)。

Eris 的质量可以通过更高精度计算得出。基于当时公认的 Dysnomia 轨道周期数值——15.774 天\(^{[24][69]}\)，Eris 的质量比 Pluto 大 27\%。使用 2011 年掩食结果，Eris 的密度为 2.52±0.07 g/cm\(^3\)\(^{[g]}\)，明显高于 Pluto，因此其主要成分必然是岩石物质\(^{[13]}\)。

内部放射性衰变加热模型表明，Eris 在地幔—核心边界处可能存在一层地下液态水海洋\(^{[70]}\)。由其卫星 Dysnomia 引起的潮汐加热也可能有助于地下海洋的持续存在\(^{[71]}\)。进一步研究表明，Eris、Pluto 和 Makemake 可能都存在活跃的地下海洋，并表现出活跃的地热活动\(^{[72]}\)。

2015 年 7 月，在 Eris 几乎被认为是太阳系中第九大直接绕 Sun 运行的已知天体近十年后，New Horizons 任务的近距离成像确定 Pluto 的体积略大于 Eris\(^{[73]}\)。因此，就体积而言，Eris 是已知直接绕 Sun 运行的天体中第十大者，而就质量而言则为第九大者\(^{[citation needed]}\)。
\begin{figure}[ht]
\centering
\includegraphics[width=6cm]{./figures/43d5ea8879c16b62.png}
\caption{} \label{fig_Eris_9}
\end{figure}
\subsection{表面与大气}
\begin{figure}[ht]
\centering
\includegraphics[width=6cm]{./figures/c6a2877412a75297.png}
\caption{与冥王星相比，Eris的红外光谱显示了两个物体之间的明显相似之处。箭头表示甲烷吸收线。} \label{fig_Eris_10}
\end{figure}
发现团队在最初识别 Eris 后，于 2005 年 1 月 25 日使用位于 Hawaii 的 8 m Gemini North Telescope 进行了光谱观测。来自该天体的红外光显示存在甲烷冰，这表明其表面可能与 Pluto 类似，当时 Pluto 是唯一已知具有表面甲烷的 TNO，而 Neptune 的卫星 Triton 其表面也含有甲烷\(^{[74]}\)。2022 年，James Webb Space Telescope（JWST）对 Eris 进行的近红外光谱观测发现其表面存在氘化甲烷冰，其丰度低于木星族彗星（如 67P/Churyumov–Gerasimenko）中的含量\(^{[75]}\)。Eris 相对较低的氘丰度表明，其甲烷并非原初来源，而可能由地下地球化学过程产生\(^{[75]}\)。JWST 还检测到了 Eris 上存在大量氮冰，推测其来源与 Eris 可能的“非原初甲烷”类似，均来自地下过程\(^{[75]}\)。估计 Eris 上氮冰的体积丰度约为甲烷的三分之一\(^{[75]}\)。

不同于 Pluto 和 Triton 略带红色且表面多样的外观，Eris 的表面几乎呈白色且均一\(^{[2]}\)。Pluto 的红色被认为是由于其表面沉积了 tholin，而这些沉积物使表面变暗，降低反照率，导致温度升高并使甲烷沉积物蒸发。相比之下，Eris 距 Sun 足够遥远，即使在反照率较低的区域，甲烷也能在表面凝结。甲烷在整个表面均匀凝结会减少任何反照率差异，并覆盖所有红色 tholin 沉积\(^{[32]}\)。这种甲烷升华与凝结循环可能在 Eris 上产生类似 Pluto 的刀刃状地形\(^{[76][75]}\)。另一种可能性是 Eris 的表面通过全球甲烷与氮冰冰川的放射性对流而得到更新，类似 Pluto 的 Sputnik Planitia\(^{[77][75]}\)。JWST 的光谱观测结果支持 Eris 表面持续更新的观点，因为在 Eris 表面未检测到乙烷这一甲烷辐解的副产物\(^{[75]}\)。

由于 Eris 的轨道遥远且偏心，其表面温度估计在约 30 至 56 K（−243.2 至 −217.2 °C；−405.7 至 −358.9 °F）之间\(^{[2]}\)。尽管 Eris 与 Sun 的距离可比 Pluto 远三倍，但当其接近 Sun 时，表面的一些冰可能会升华形成大气。由于甲烷和氮都高度挥发，它们的存在表明，Eris 要么始终位于太阳系遥远且寒冷的区域，使甲烷与氮冰得以长期存在，要么天体本身具有内部机制以补充从大气中逃逸的气体\(^{[76]}\)。这与另一已发现的 TNO——Haumea 的观测结果形成对比，后者显示存在水冰但未检测到甲烷\(^{[78]}\)。
\subsection{自转}  
由于 Eris 表面均匀，其在自转过程中亮度变化极小，从而使其自转周期的测量变得困难\(^{[79][16]}\)。对 Eris 亮度的长期精确监测表明，它与其卫星 Dysnomia 发生了潮汐锁定，其自转周期与卫星轨道周期 15.78 个 Earth 日同步\(^{[16]}\)。Dysnomia 也与 Eris 潮汐锁定，使 Eris–Dysnomia 系统成为继 Pluto 与 Charon 之后第二个已知的“双同步自转”系统。此前由于缺乏对 Eris 自转的长期观测，得到的自转周期测量值非常不确定，范围从数十小时到数天不等\(^{[79][80][81]}\)。Eris 的自转轴倾角尚未测定\(^{[15]}\)，但合理推测其与 Dysnomia 的轨道倾角一致，即相对于黄道面约为 78°\(^{[17]}\)。若情况如此，Eris 北半球的大部分区域将被 Sun 光照亮，其中约 30\% 的半球在 2018 年将处于持续照明状态\(^{[17]}\)。
\subsection{卫星}
\begin{figure}[ht]
\centering
\includegraphics[width=6cm]{./figures/557e31975ded2068.png}
\caption{艺术构想图：Eris 及其暗色卫星 Dysnomia} \label{fig_Eris_11}
\end{figure}
2005 年，Hawaii 的 Keck 望远镜自适应光学团队使用新启用的激光导星自适应光学系统，对四颗最明亮的 TNO（Pluto、Makemake、Haumea 和 Eris）进行了观测\(^{[82]}\)。9 月 10 日拍摄的图像显示，Eris 周围存在一颗轨道卫星。按照当时 Eris 已被广泛使用的昵称 “Xena”，Brown 团队将这颗卫星昵称为 “Gabrielle”，取自电视系列剧中战士公主（Xena）的伙伴名字。当 Eris 获得 IAU 正式命名后，该卫星被命名为 Dysnomia，取自希腊“不法之女神”，即 Eris 的女儿。Brown 表示，他选择这一名称是因为其发音与其妻子 Diane 相似。该名称还含蓄保留了与 Eris 早期昵称 Xena 的联系（Xena 由电视演员 Lucy Lawless 扮演），尽管这一联系并非故意为之\(^{[83]}\)。
\begin{figure}[ht]
\centering
\includegraphics[width=10cm]{./figures/0684e3bc9147f253.png}
\caption{} \label{fig_Eris_12}
\end{figure}
\subsection{探测}  
在 2020 年 5 月，Eris 被正在远离太阳系的 New Horizons 探测器从远距离观测，这属于其 2015 年成功飞掠 Pluto 之后扩展任务的一部分\(^{[18]}\)。尽管当时 Eris 与 New Horizons 的距离（112 AU）大于其与 Earth 的距离（96 AU），但由于探测器处于 Kuiper 带内部这一独特观测位置，可以获得从 Earth 无法实现的高相位角观测，从而确定 Eris 表面的光散射特性和相位曲线行为\(^{[18]}\)。

在 2010 年代，曾进行多项关于后续 Kuiper 带探测任务的研究，其中将 Eris 作为候选目标进行评估\(^{[85]}\)。计算显示，若采用 Jupiter 引力助推，对于分别在 2032 年 4 月 3 日或 2044 年 4 月 7 日发射的任务，飞掠 Eris 所需时间为 24.66 年。当探测器抵达时，Eris 与 Sun 的距离将分别为 92.03 AU 或 90.19 AU\(^{[86]}\)。
\subsection{另见}
\begin{itemize}
\item 天体命名
\item 清除邻近的小天体
\item 国际天文学联合会（IAU）
\item 假设的海王星外行星
\end{itemize}
\subsection{参考文献}
\begin{enumerate}
\item Staff. Discovery Circumstances: Numbered Minor Planets. IAU: Minor Planet Center. 2007-05-01 [2007-05-05]. （原始内容存档于2016-05-04）.
\item Brown, Mike. The discovery of 2003 UB313 Eris, the largest known dwarf planet. 2006 [2007-05-03]. （原始内容存档于2014-09-10）.
\item JPL Small-Body Database Browser: 136199 Eris (2003 UB313). 2009-11-20 [2009-01-21]. （原始内容存档于2019-01-09）.
\item List Of Centaurs and Scattered-Disk Objects. Minor Planet Center. [2008-09-10]. （原始内容存档于2011-07-25）.
\item Buie, Marc. Orbit Fit and Astrometric record for 136199. Deep Ecliptic Survey. 2007-11-06 [2007-12-08]. （原始内容存档于2014-09-26）.
\item Asteroid Observing Services （页面存档备份，存于互联网档案馆）
\item Sicardy, B.; Ortiz, J. L.; Assafin, M.; Jehin, E.; Maury, A.; Lellouch, E.; Gil-Hutton, R.; Braga-Ribas, F.; Colas, F.; Widemann. Size, density, albedo and atmosphere limit of dwarf planet Eris from a stellar occultation (PDF). European Planetary Science Congress Abstracts. 2011, 6: 137 [September 14, 2011]. Bibcode:2011epsc.conf..137S. （原始内容存档 (PDF)于October 18, 2011）.
\item Beatty, Kelly. Former 'tenth planet' may be smaller than Pluto. NewScientist.com. Sky and Telescope. November 2010 [October 17, 2011]. （原始内容存档于February 23, 2012）.
\item 掩星觀測顯示鬩神星和冥王星大小幾乎相等. PanSci 泛科学. 2011-11-03 [2017-03-09]. （原始内容存档于2014-03-25） （中文（台湾））.
\item Michael E. Brown and Emily L. Schaller. The Mass of Dwarf Planet Eris (abstract page). Science. 2007, 316 (5831): 1585 [2011-02-01]. PMID 17569855. doi:10.1126/science.1139415. （原始内容存档于2010-06-26）.
\item Kelly Beatty. Eris Gets Dwarfed (Is Pluto Bigger?). Sky & Telescope (News Blog). 2010-11-07 [2010-11-07]. （原始内容存档于2011-08-30）.
\item Snodgrass, Carry, Dumas, Hainaut. Characterisation of candidate members of (136108) Haumea's family (abstract). The Astrophysical Journal. 2009-12-16. arXiv:0912.3171.
\item AstDys (136199) Eris Ephemerides. Department of Mathematics, University of Pisa, Italy. [2009-03-16]. （原始内容存档于2011-06-04）.
\item F. Bertoldi, W. Altenhoff, A. Weiss, K.M. Menten, C. Thum. The trans-neptunian object UB313 is larger than Pluto. Nature. 2006-02, 439 (7076): 563–564 [2018-04-02]. ISSN 1476-4687. doi:10.1038/nature04494. （原始内容存档于2019-07-13） （英语）.
\item “阋”，《说文解字》卷三：“恒讼也。《诗》云：‘兄弟阋于墙。’从斗从儿。儿，善讼者也。许激切。”繁：“鬩”，简：“阋”，拼音：xì，注音：ㄒㄧˋ，粤拼：jik1
\item JPL/NASA. What is a Dwarf Planet?. Jet Propulsion Laboratory. 2015-04-22 [2022-01-19]. （原始内容存档于2021-12-08）.
\item How Big Is Pluto? New Horizons Settles Decades-Long Debate. www.nasa.gov. 2015 [2015-07-14]. （原始内容存档于2017-07-01）.
\item Mike Brown. The discovery of 2003 UB313 Eris, the 10th planet largest known dwarf planet. Caltech. 2006 [2010-01-05]. （原始内容存档于2012-05-20）.
\item The IAU draft definition of "planet" and "plutons" (新闻稿). IAU. 2006-08-16 [2006-08-16]. （原始内容存档于2006-08-20）.
\item Mike Brown. The shadowy hand of Eris. Mike Brown's Planets. 2010 [2010-11-07]. （原始内容存档于2010-11-11）.
\item Mike Brown. How big is Pluto, anyway?. Mike Brown's Planets. 2010-11-22 [2010-11-23]. （原始内容存档于2015-10-31）. (Franck Marchis on 2010-11-08)
\item Brown, Mike. How big is Pluto, anyway?. Mike Brown's Planets. 2010-11-22 [2010-11-23]. （原始内容存档于2015-10-31）. (Franck Marchis on 2010-11-08)
\item Thomas H. Maugh II and John Johnson Jr. His Stellar Discovery Is Eclipsed. Los Angeles Times. 2005-10-16 [2008-07-14]. （原始内容存档于2017-02-21）.
\item KLC. 矮行星Eris中譯名正式命名為「鬩神星」. tamweb.tam.gov.tw. 台北天文馆. 2007-07-10 [2017-03-09]. （原始内容存档于2013-07-29） （中文（台湾））.
\item 方舟子. 方舟子：欺世盗名的八卦宇宙论. 人民网. 2005-08-01 [2013-03-30]. （原始内容存档于2013-05-09） （中文（中国大陆））.
\item Gemini Observatory Shows That"10th Planet" Has a Pluto-Like Surface. Gemini Observatory. 2005 [2007-05-03]. （原始内容存档于2007-03-11）.
\item M. E. Brown, C. A.Trujillo, D. L. Rabinowitz. Discovery of a Planetary-sized Object in the Scattered Kuiper Belt. The Astrophysical Journal (abstract page). 2005, 635 (1): L97–L100. Bibcode:2005ApJ...635L..97B. arXiv:astro-ph/0508633 可免费查阅. doi:10.1086/499336.
\item J. Licandro, W. M. Grundy, N. Pinilla-Alonso, P. Leisy. Visible spectroscopy of 2003 UB313: evidence for N2 ice on the surface of the largest TNO (PDF). Astronomy and Astrophysics. 2006, 458 (1): L5–L8 [2011-05-03]. Bibcode:2006A&A...458L...5L. arXiv:astro-ph/0608044 可免费查阅. doi:10.1051/0004-6361:20066028. （原始内容存档 (PDF)于2008-11-21）.
\item M. E. Brown, M. A. van Dam, A. H. Bouchez, D. LeMignant, C. A. Trujillo, R. Campbell, J. Chin, Conrad A., S. Hartman, E. Johansson, R. Lafon, D. L. Rabinowitz, P. Stomski, D. Summers, P. L. Wizinowich. Satellites of the largest Kuiper belt objects. The Astrophysical Journal (abstract page). 2006, 639 (1): L43–L46. Bibcode:2006ApJ...639L..43B. arXiv:astro-ph/0510029 可免费查阅. doi:10.1086/501524.
\end{enumerate}
\subsection{外部链接}
\begin{itemize}
\item Michael Brown 关于 Eris 和 Dysnomia 的网页  
\item MPC 关于 (136199) Eris 的数据库条目  
\item NASA JPL 的 3D 轨道可视化  
\item 2003 UB313 轨道模拟  
\item Keck 天文台关于 Dysnomia 发现的页面  
\item AstDyS-2（小行星动力学网站）中的 Eris（矮行星）  
\item 历表 · 观测预测 · 轨道信息 · 本征要素 · 观测信息  
\item JPL 小天体数据库中的 Eris（可在 Wikidata 编辑）  
\item 近距离接近 · 发现 · 历表 · 轨道查看器 · 轨道参数 · 物理参数
\end{itemize}












